\section{Reproducibility Details}\label{app:reproducibility}

\subsection{Source preparation and refinement settings}\label{app:preprocessing}

The AMASS source collections used by the preparation pipeline are ACCAD,
BioMotionLab\_NTroje, BMLhandball, BMLmovi, CMU, DFaust\_67, EKUT,
Eyes\_Japan\_Dataset, HumanEva, KIT, MPI\_HDM05, MPI\_Limits, MPI\_mosh,
SFU, SSM\_synced, TCD\_handMocap, TotalCapture, and Transitions\_mocap.
Source preparation operates at 20\,Hz; the final motion libraries use
30\,Hz. Joint ordering and quaternion conventions must follow the matching
asset and MotionLib conversion rather than be inferred from source array
positions.

\begin{table}[htbp]
\centering
\small
\caption{Default settings of the current reference-refinement procedure.
Frame counts refer to the 30\,Hz MotionLib representation. These are
kinematic-processing settings, not physical constraints learned by the policy.}\label{tab:refinement_parameters}
\begin{tabularx}{\linewidth}{@{}lY@{}}
\toprule
Operation & Setting \\
\midrule
Contact candidates & Left/right ankle and toe bodies \\
Contact detection thresholds & Speed 0.15\,m/s; height 0.10\,m \\
Contact cleanup & Close gaps up to 2 frames; remove runs shorter than 3 frames \\
Rotation filter & Savitzky--Golay, 9-frame window, polynomial order 3 \\
Rotation blend / bound & Strength 0.5; maximum geodesic correction $8^\circ$ \\
Contact-transition guard & Separate segments; 2-frame transition taper \\
Root XY correction & Median stance anchors; least-squares common translation \\
Root correction regularization & 21-frame smoothing window; 0.10\,m displacement cap \\
Ground clearance & 0.005\,m collision-geometry margin; 5-frame lift smoothing \\
Final reconstruction & Per-asset FK and per-motion velocity recomputation \\
\bottomrule
\end{tabularx}
\end{table}

\subsection{PD gains and limits}\label{app:control}

Table~\ref{tab:pd_parameters} gives the base values configured for the
SMPL morphology robot. For joint $d$ belonging to rigid body $j$, the
simulator applies the multiplier
\begin{equation}
    \rho_{b,j} = \frac{m_{b,j}}{m_{\mathrm{ref},j}}
\end{equation}
to stiffness, damping, and effort limit. The velocity limit is not mass-scaled.
All relevant asset masses, joint ranges, collision settings, and the reference
asset must accompany a reproducible physical interpretation.

\begin{table}[htbp]
\centering
\small
\caption{Base PD and actuation values before per-body mass scaling.
Stiffness and damping use the simulator's rotational PD units
(Nm/rad and Nm\,s/rad); effort limits are in Nm and velocity limits in rad/s.}\label{tab:pd_parameters}
\begin{tabular}{lrrrr}
\toprule
Joint family & Stiffness & Damping & Effort limit & Velocity limit \\
\midrule
Hip, knee, ankle & 800 & 80 & 500 & 100 \\
Toe & 500 & 50 & 500 & 100 \\
Torso, spine, chest & 1,000 & 100 & 500 & 100 \\
Neck, head, thorax, shoulder, elbow & 500 & 50 & 500 & 100 \\
Wrist, hand & 300 & 30 & 500 & 100 \\
\bottomrule
\end{tabular}
\end{table}

\subsection{Windowed jerk convention}\label{app:smoothness}

The logged smoothness score is computed from simulated global body positions
using consecutive finite differences with control interval
$\Delta t=1/30$\,s. A nominal 0.4\,s window contains
$W=\max(4,\lfloor0.4/\Delta t\rfloor)$ frames, giving $W=12$ at the
configured rate and duration $T_w=(W-1)\Delta t$. For body $j$ in window $w$,
let $L_{w,j}$ be the sum of speed times $\Delta t$ and let
$\dddot{p}_{w,j}$ denote the third finite difference. The implementation uses
\begin{equation}
    J_{w,j} =
    \frac{T_w^5
      \sum_{t\in w}\|\dddot{p}_{t,j}\|_2^2\,\Delta t}
    {\max(L_{w,j},\epsilon)^2+\epsilon},
    \qquad \epsilon=0.1,
    \label{eq:jerk}
\end{equation}
where the sum includes the valid third-difference samples in the window.
Window scores are averaged per body and then across bodies and motions.
High-jerk percentage counts windows for which
$\max_j J_{w,j}>6{,}500$.

Equation~\eqref{eq:jerk} documents the existing numerical convention, including
its regularizers; it is not presented as a universally dimensionless or
sampling-rate-invariant biomechanical measure. Comparisons must keep the
sample rate, window, units, regularizers, and aggregation fixed. Reference
refinement reports also include raw jerk in m/s$^3$; those values are a
different metric and cannot be compared numerically to this logged score.

\subsection{Split and experiment provenance}\label{app:provenance}

The split metadata is stored in
\path{data/splits/hhi_stage2_v1/split_metadata.json}, with schema version 1,
split seed 42, and split identifier
\begin{quote}
\footnotesize\ttfamily
4c1c9d592e3dbeef15abe18cb50a3141820735849414ecc2d8959e9f38502263
\end{quote}
It records the source-manifest hash, training/validation/test manifest and
ID-list hashes, difficulty information, and the forced training reservation.
The generator uses 100 difficulty strata; the 150 requested development
identities expand to 286 retained identities when existing mirrored partners
are included. Checkpoint/run split provenance refers to these artifacts,
not just a verbal 90/5/5 ratio.

The documented HUMOS configuration selects
\path{logs/humos/q6zbv2tu/checkpoints/latest-epoch=1599.ckpt}. The immutable
HUMOS checkpoint checksum and the SMPLSim / asset-generation revision hashes
are to be archived with the dataset release; a checkpoint path identifies the
local workflow but is not a substitute for a checksum.

The principal implementation entry points are:
\begin{itemize}
    \item Reference refinement:
    \path{tools/refine_humos_motion.py}; streamed per-clip processing:
    \path{tools/refine_humos_r2_per_clip.py}.
    \item Selected full-scale experiment:
    \path{examples/experiments/mimic/mlp_wide_stage2_discover_attention_slot_type.py},
    inheriting the common attention experiment and enabling slot/type embeddings.
    \item Corrected development comparison:
    \path{tools/train_150_evalfix_ablation.sh} (cases A--E) and
    \path{tools/train_150_evalfix_unrefined.sh} (cases F--H), seed 0.
    \item Motion residency and evaluation:
    \path{protomotions/components/global_clip_pool.py} and
    \path{protomotions/agents/evaluators/mimic_evaluator.py}.
    \item Standalone split evaluation:
    \path{evaluate_validation_split.py} and \path{evaluate_test_split.py}
    under \path{protomotions/}. Each output records the checkpoint, split,
    reference version, and shape panel; the test evaluator additionally
    checks the test-manifest hash, count, and disjointness before running.
\end{itemize}

For each final experiment, retain the resolved configuration, code revision,
model/checkpoint checksum, body coefficient and asset manifests, split hashes,
all resume segments, frame count, GPU model/count, and elapsed GPU-hours.
For each evaluation, retain per-pair clip and asset IDs, panel selection,
success/failure, valid horizon, and component errors in a local analysis file.
The main full-scale W\&B identifier is \texttt{3mvgad6f}; the corrected
small-scale group is \texttt{hhi\_150\_evalfix\_comparison}.
The evaluation/sampling correction is recorded in commit \texttt{3685247}.

\pending{Two items remain before a finalized benchmark report: the immutable
HUMOS/SMPLSim checksums (to be archived with the dataset release) and the
release of the retained per-motion evaluation records for the validation and
test splits.}
